\chapter{Khởi Động Bằng Ba Mươi Giây}
\chapterintro{Nhanh, gọn, tỉnh lớp}{Không phải lúc nào mở tiết cũng cần đủ 5 phút; có những ngày chỉ cần 30 giây đúng chỗ là lớp vào nhịp.}

\section{Ba tiếng gọi đồng loạt}
\begin{khoidongbox}{Hoạt động khởi động}
Mở bằng ba tiếng rất ngắn: \textit{``Ngẩng lên nào''}, rồi dừng hai giây để cả lớp tự gom ánh mắt về bảng.
\end{khoidongbox}
\begin{visaohieuqua}
Lớp phản ứng tốt với tín hiệu ngắn và rõ hơn là một đoạn nhắc trật tự dài. Ba tiếng này đủ nhẹ mà vẫn có lực.
\end{visaohieuqua}
\begin{cachlam5phut}
Dùng trong 30 giây đầu, sau đó nối ngay bằng một câu mở bài hoặc mục tiêu tiết học. Không lặp quá hai lần.
\end{cachlam5phut}
\ngancach

\section{Chọn một con số để vào tiết}
\begin{khoidongbox}{Hoạt động khởi động}
Hỏi: \textit{``Nếu chấm trạng thái của mình từ 1 đến 5, em đang ở mức mấy?''}
\end{khoidongbox}
\begin{visaohieuqua}
Con số giúp học sinh phản hồi cực nhanh. Thầy cô cũng nhìn được nhiệt độ lớp ngay từ đầu.
\end{visaohieuqua}
\begin{cachlam5phut}
Cho giơ tay hoặc giơ ngón tay, rồi chốt bằng câu: \textit{``Rồi, mình kéo thêm một nấc nữa.''}
\end{cachlam5phut}
\ngancach

\section{Một câu dự báo cho bài học}
\begin{khoidongbox}{Hoạt động khởi động}
Mời lớp đoán rất nhanh: \textit{``Hôm nay bài này sẽ dễ ở chỗ nào, khó ở chỗ nào?''}
\end{khoidongbox}
\begin{visaohieuqua}
Khi được dự báo trước, học sinh bước vào bài với tâm thế chủ động hơn. Các em không còn chỉ ngồi chờ nghe.
\end{visaohieuqua}
\begin{cachlam5phut}
Lấy 2 ý ngắn rồi viết một từ khóa lên bảng. Từ đó chuyển thẳng sang phần chính.
\end{cachlam5phut}
\ngancach

\section{Giơ tay theo mức tỉnh}
\begin{khoidongbox}{Hoạt động khởi động}
Nói: \textit{``Ai thấy mình đã sẵn sàng thì giơ một tay, ai cần cứu nhẹ thì giơ hai tay.''}
\end{khoidongbox}
\begin{visaohieuqua}
Cả lớp cùng tham gia trong vài giây và không khí bật lên rất nhanh. Đây là kiểu khởi động không cần phát biểu dài.
\end{visaohieuqua}
\begin{cachlam5phut}
Sau khi lớp phản hồi, chọn luôn một động tác cứu nhịp ngắn như hít sâu, nói một từ, hoặc trả lời một câu rất dễ.
\end{cachlam5phut}
\ngancach

\section{Nói nối một từ}
\begin{khoidongbox}{Hoạt động khởi động}
Mời ba học sinh nối tiếp nhau, mỗi em nói đúng một từ để tạo thành câu mở đầu cho tiết học.
\end{khoidongbox}
\begin{visaohieuqua}
Một từ thì ai cũng dám nói. Khi ba từ ghép lại, lớp thường bật cười và tỉnh rất nhanh.
\end{visaohieuqua}
\begin{cachlam5phut}
Chọn ngẫu nhiên 3 em, chốt lại câu vui nhất và dùng nó làm cầu nối vào bài.
\end{cachlam5phut}